-----5.6.3.6: Best Settings-----


One of the experimental objectives was to
determine "best" settings for (X1,X2,X3,X4,X5) so
as to achieve Y = 30, 60, and 90 inch distances.
To solve this problem, we utilize 2 approaches:

   1. Graphical     : DEX Contour Plot
   2. Quantititative: Sorted Data Table

The virtue of these approaches is that they both
depend on the data itself, and may be done
independently of the previously-discussed modeling
and any of the deficiencies therein.

-----Section 1: DEX Contour Plot-----

The only piece of information needed from the
prior analysis is the top 2 entrees in the ranked
list of factors:

      1. X4  (arm length)      (effect =  40.3 inches)
      2. X3  (number of bands) (effect =  35.9 inches)
      3. X1  (band height)     (effect =  27.0 inches)
      4. X5  (start point)     (effect =  24.0 inches)
      5. X2  (stop angle)      (effect = -22.2 inches)
      6. X3*X4 interaction     (effect =  15.2 inches)
      7. X1*X4 interaction     (effect =   9.4 inches)
      8. X1*X3 interaction     (effect =   9.3 inches)

namely, that X4 (arm length) and X3 (number of
bands) are the two most important factors.  We
thus choose to use these 2 factors as the axes on
the contour plot:

      X4 (arm length)     : horizontally
      X3 (number of bands): vertically

The values at the 4 vertices of this data square
are gotten by simple averages of the data points.
Thus, for example, at the vertex (X4 = +1, X3 =
-1) the 4 data points from the original data table
which have x4 = +1 and X3 = -1 are 33, 85, 45, and
36.5.  These 4 values average to 45.5 which
becomes the value at that vertex.  We do similarly
for the other 3 vertices.  As for the
pseudo-center point at X3 = -1, the data is 45 and
37.5 with the average being 41.25.  At the
pseudo-center point at X3 = +1, the data is 99 and
84.5 thus yielding an average of 91.75.  The
original 20 data points are thus are reduced to 6
mean values.  Based on these mean values, the DEX
contour plot (with contour lines at 30, 40, 50,
..., 90) is as follows:

   <dex contour plot>

Conclusions:
From the contour plot, we note the following:
   1. Curvature & Interactions: The curvature in
      the contour lines indicates a strong
      interaction effect between X3 and X4.
      Referring back to the Yates table ranking,
      we note that the X3*X4 interaction was in
      fact the largest of all of the intereactions
      (with a magnitude of 15.2), and so the above
      graphical curvilinearity is consistent with
      prior quantitative results.  Note conversely
      that if the contour curves had been linear
      then that would imply a planar model in X3
      and X4 with the additional implication that
      the X3*x4 cross-product effect is near zero.

   2. Cross-Validation: This is the statistical
      procedure whereby a part (most) of the data
      is used for estimation and inference, and a
      second (smaller) part of the data is set
      aside and is used for validation and the
      confirmation.  Such was used here: contour
      curves were (by construction) based on the
      16 edge data points--the 4 pseudo-center
      points were NOT used in the computation of
      the contour curves.  It is known (small
      residual standard deviation) that the
      contour curves predict well on the edges.
      Can we confidently use the model for
      prediction everywhere else?  Anywhere else?
      particular, does it predict well interior to
      the cube?  To answer this, let us compare
      the contour curves to the known data
      averages at the 2 pseudo-center points (at
      X3 = -1 (41.25) and at X3 = +1 (91.75)).  If
      there is a good match, then that gives us
      confidence in using the edge-based model for
      interpolation; if not, then the use of the
      edge-based model is severely restricted.  In
      this case, we note (eyeball) that at the X3
      = -1 pseudo-center point, the average is
      41.25 and the contour-lined based prediction
      is about 37 and so the error is
      approximately 4 inches.  For the X3 = +1
      pseudo-center point, the average is 91.75
      but the contour-lined based prediction is
      about 75 and so the error in prediction is
      an overly-large 16 inches.  Hence in this
      case, the contour lines predict well in the
      vicinity of the X3 = -1 pseudo-center point,
      but poorly at the X3 = +1 pseudo-center
      point.  This cross-validation is discuraging
      and hence implies that the model may NOT be
      used freely for interpolation, and hence
      interpolatory best settings based on the
      model are not to be trusted.

   3. Discrete Factor: Note that the factor X3
      (number of bands) is intrinsically discrete
      (the number of bands is either 1 or 2 and
      one could not have, e.g., 1.7 bands).  Given
      this, then one may have argue that the
      contour plots are meaningless.  In a very
      real sense that is correct; on the other
      hand, the above 2 conclusions are both
      correct and the contour plots--discreteness
      notwithstanding--helped arrive at these 2
      conclusions.  In preparation for the
      upcoming question as to what are the best
      settings for (X1,X2,X3,X4,X5) for Y = 30,
      60, and 90, we make use of the discreteness
      of X3 and X3's relatively high importance
      (the second most important factor) to split
      the data (based on X3= -1 and X3 = +1) and
      generate the following 2 contour plots of X1
      (the third most important factor) versus X4
      (the most important factor)

         <2 subset contour plots>

   4. Conclusions from the X3 = -1 Contour Plot:
      For X3 = -1 (that is, for the number of bands = 1),
         1. the curves are relatively linear
            (implying a relatively small X1*X4 interaction);
         2. the response curves are relatively small
            (30 to 60);
         3. the curves fit the center point relatively well
            (actual = 41.25; predicted = 37).

   5. Conclusions from the X3 = +1 Contour Plot:
      For X3 = +1 (that is, for the number of bands = 2),
         1. the curves are curvilinear
            (implying a large X1*X4 interaction);
         2. the response curves are relatively large
            (40 to 90);
         3. the curves fit the center point poorly
            (actual = 91.75; predicted = 73).

   6. Best Settings for Y = 30:
      Choose the level of the splitting factor
      first; in this case, since X3 = -1 is more
      centered on lower values of Y than X3 = +1,
      and since X3 = -1 has better interpolatory
      properties than X3 = +1, then choose X3 =
      -1.  Given that, any point along the Y = 30
      curve will work.  We choose the intersection
      point of the Y = 30 curve and the bottom of
      the box (X1 = -1); thus yielding (X4 = +.8,
      X1 = -1).  This is also near a vertex point
      with average 31.625.  We hence have 3 out of
      the 5 settings:

         (X1 = -1, X2 = ?, X3 = -1, X4 = +0.8, X5 = ?)

      The remaining 2 settings (for X2 and X5)
      will be derived later.

   7. Best Settings for Y = 60:
      Since X3 = -1 has better interpolatory
      properties than X3 = +1, then choose X3 =
      -1.  As before, any point along the Y = 60
      curve will work.  Note that Y = 60 comes
      very close to the (X4 = +1, X1 = +1) vertex
      which has an (average) value of 60.75--very
      close to 60.  Unfortunately, the 2 raw data
      values going into the 60.75 are quite
      disparate: 36.5 (from ++-+-) and 85.0 (from
      +--++).  We thus see that the settings for
      X2 and X5 make a huge difference in the
      response: (X2 = +1, X5 = -1) yielding 36.5
      while (X2 = -1, X5 = +1) yielding 85.0.  We
      will decide on these 2 value later, but for
      now, as before, we have

         (X1 = -1, X2 = ?, X3 = -1, X4 = +0.8, X5 = ?)

   8. Best Settings for Y = 90:
      Since the X3 = -1 contour plot has no values
      in the vicinity of Y = 90 except for distant
      extrapolation, we then resort to the X3 = +1
      plot.  On this plot, we could make use of
      the Y = 90 curve, but since the X3 = +1 case
      has such poor interpolatory properties, then
      it makes better sense to set aside the
      biased curves and rely more heavily on the
      observed data.  Note, for example, that the
      center point for the X3 = +1 contour plot
      has an average value of 91.75--very close to
      90.  Note that the data which led to this
      center point average of 91.75 is from
      (00100) replicates: 99.0 and 84.5.  These 2
      data points--even while being collected
      under identical conditions--still differed
      by almost 15 inches (with a standard
      deviation of 10.25 inches).  This is large
      but is as good as the data can do--having 2
      values under replicate conditions
      symmetrically spanning the desired Y = 90.
      Thus in conclusion we choose for our best
      settings for Y = 90 to be

            (X1 = 0, X2 = 0, X3 = +1, X4 = 0, X5 = 0)

      This is primarily a data-based estimate for
      the best settings, but given that the model
      performs poorly and yields biased estimates
      for this X3 = +1 case, such a data-based
      estimate is the preferred choice.

-----Section 2: Sorted Data-----

Up to this point, there has been a significant
effort in obtaining best settings based on
graphical tools in conjuntion with best-fit
modeling.  The last step is a quantitative one and
involves a simple ordering of the data (smallest
to largest).  In the absence of any other
techniques, such an ordering is quite useful in
terms of determining best settings, especially if
the observed data happens to fall in the vicinity
of the desired target response values, namely, Y =
30, 60, and 90.

The sorted data (carrying along the settings of X1
to X5 is as follows:


                                                 RUN  CENT
   Y      X1  X2  X3  X4  X5 SEQ POINT
----------------------------------------
  8.00    -1   1  -1  -1  -1  10     0
 28.00    -1  -1  -1  -1   1   1     0
 28.25     1   1  -1  -1   1   8     0
 28.50     1   1   1  -1  -1  14     0

 33.00    -1  -1  -1   1  -1  12     0
 33.50    -1  -1   1  -1  -1  15     0
 35.00     1  -1  -1  -1  -1   6     0
 36.00    -1   1   1  -1   1  16     0
 36.50     1   1  -1   1  -1  11     0
 37.50     0   0  -1   0   0  19     2
 45.00    -1   1  -1   1   1  18     0
 45.00     0   0  -1   0   0   7     2
 45.00    -1   1   1   1  -1   5     0

 84.00     1  -1   1  -1   1  17     0
 84.50     0   0   1   0   0  13     1
 85.00     1  -1  -1   1   1   9     0

 99.00     0   0   1   0   0   2     1
106.00    -1  -1   1   1   1  20     0
126.50     1  -1   1   1  -1   4     0
126.50     1   1   1   1   1   3     0


Conclusions:
   1. Bimodlality of the Data: Note that the data
      naturally splits itself nto 2 regions: 45
      and below (13 values), and 84 and above (7
      values).

   2. Important Factors:
      X3: Of the 13 low values, 9 have X3 = -1.
          Of the 7 highest values, 6 have X3 = +1
          This kind of strong reversal implies X3 is
          the most important of the factors.

      X4: Of the 13 lowest values, 7 have X4 = -1.
          Of the 7 highest values, only 1 has X4 = +1.
          The 4 lowest values have X4 = -1.
          The 3 highest values have X4 = +1.
          X4 is also important (but next so).

      X1: Of the 13 lowest values, 7 have X1 = -1.
          Of the 7 highest values, only 1 has X1 = +1.
          The 2 lowest values have X1 = -1.
          The 2 highest values have X1 = +1.
          X1 is next in importance.

      X5: Of the 13 lowest values, 7 have X5 = -1.
          Of the 7 highest values, only 1 has X5 = +1.
          X5 is next in importance.

      X2: The distribution of -1's and +1's
          is fairly even.
          X2 is the least important of the 5 factors.

   3. Best Settings for Y = 30:
      Based only on the above sorted data, and
      considering the large replication standard
      deviation of the data, a good choice for Y =
      30 would be the closest data point, namely,
      28.5 which comes via

         (X1 = +1, X2 = +1, X3 = +1, X4 = -1, X5 = -1)

      Although this point has the smallest error
      (28.5 - 30 = -1.5 inches), any of the close
      points may also be used:

            28.5     + + + - -
            28.25    + + - - +
            28.0     - - - - +
            33.0     - - - + -
            33.5     - - + - -

   4. Best Settings for Y = 60:
      Based only on the above ordered data, no
      setting yields a response anywhere in the
      vicinity of Y = 60 and so no data-based
      best setting is possible.  In fact, there is no
      empirical proof that ANY (possible) settings
      of the catapult system will yield Y = 60.
      We generate Y = 60 predicted values based on
      the fact that we are assuming the model to
      be of correct form and to be valid over the entire
      region.  The first assumption is
      questionable and the second is false due to
      poor interpolatory prediction (especially
      for X3 = +1).

   5. Best Settings for Y = 90:
      Based only on the above sorted data, and
      considering the large replication standard
      deviation of the data, a good choice for Y =
      90 would be the closest data point, namely,
      85 which comes via

         (X1 = +1, X2 = -1, X3 = -1, x5 = +1, x6 = +1).

      Although this point has the smallest error:
      85 - 90 = -5 any of the close points may
      also be used:

            85.0     + - - + +
            84.5     0 0 + 0 0
            84.0     + - + - +

-----Section 3: Recommendations-----

1. Best Settings for Y = 30:
   If simplicity is preferred, use the "Sorted
   Data" solution (+ + + - -) which yielded a
   close to 30 data value (28.5).  This solution
   has the advantage that it is model-free.

   A better solution (but involves collecting more
   data) is to use the "DEX Contour Plot" solution
   of (- ? - + ?).  The values of X2 and X5 are
   unknown.  There are 4 cases for X2 and X5:

      - - - + -   data = 33.0
      - + - + +   data = 45.0
      - + - + -   no data
      - - - + +   no data

   Make 2 additional runs to collect data at -+-+-
   and ---++ and then draw the X5 versus X2 DEX
   contour plot and use any value of X2 and X5
   along the Y = 30 contour curve.

2. Best Settings for Y = 60:
   Since there is no data anywhere close to Y =
   60, there is no "Sorted Data" recommendation.

   The recommended solution involves collecting
   more data.  data) is to use the "DEX Contour
   Plot" solution of (+ ? - + ?).  The values
   of X2 and X5 are unknown.  There are 4 cases
   for X2 and X5:

      + + - + -   data = 36.5
      + - - + +   data = 85.0
      + - - + -   no data
      + + - + +   no data

   Make 2 additional runs to collect data at
   +--+- and ++-++ and then draw the X5 versus
   X2 DEX contour plot and use any value of X2
   and X5 along the Y = 60 contour curve.
   Further, for this noisy region it would be
   preferable to collect a center point at
   +0-+0, and even better still to replicate
   that center point to determine a local
   estimate of noise.

3. Best Settings for Y = 90:
   If simplicity is preferred, use the "Sorted Data"
   solution (0 0 + 0 0) which yielded an average close to
   the desired 90 value (91.5).  This solution has the
   advantage that it is model-free.

   Since the 2 data values (84.5 and 99) that went into that
   replicate average are quite disparate, then sampling
   anywhere else in the vicinity will probably be for
   naught.  Hence in this case, the simplest solution (at
   the replicated center point for X3 = +1) is also the
   recommended best solution.

-----5.6.3.7: Conclusion and Summary-----

Important factors and best settings have been derived via 
the usual analysis approach.  Important factors were X3, X4,
X1, X5 and weakly X2.  Some interactions (especially X3*X4) 
were also important.  Best settings included: 

   For Y = 30: (xx) or augment the (yy) data
   For Y = 60: augment the (xx) data
   For Y = 90: (0 0 + 0 0)

Anomalies existed in the data.  Plots revealed higher 
variability for X3 = +1 than for X3 = -1 implying a serious 
violation of the constant variation assumption for least 
squaes fitting.  Cross-validation revealed poor-quality fits
for the X3 = +1 data--especially at the X3 = +1 
pseudo-center point; hence interpolation was suspect.  
Corrective action is thus suggested: 

   1. Transformations of Y (to correct non-constant variation)
   2. Subset Analyses (to correct poor-fitting regions)

The following section carries out such a corrected analysis. 


-----5.6.3.7: Corrected Analysis-----

